\chapter{Svar på övningarna}
\label{ch:svar}

Här finns svaren på samtliga övningar. Läs dem efter att ni försökt själva --- ett svar man läst
är inte samma sak som ett svar man kommit fram till, och det är den skillnaden övningarna finns
för.

Där ett svar bygger på en uträkning står uträkningen med, inte bara resultatet.

\section{Bussen}

\solution{ex:1:1}{Kablage}
\ans{a} $8 \cdot 7 / 2 = 28$ förbindelser.
\ans{b} 8 anslutningar, en per enhet.
\ans{c} Vi söker $n$ där $n(n-1)/2 = 2n$, alltså $n - 1 = 4$ och $n = 5$. Redan vid fem enheter
är punkt-till-punkt dubbelt så mycket koppar, och avståndet växer sedan snabbt.

\solution{ex:1:2}{Bussens nivå}
\ans{a} \dom. Bussen är dominant så snart minst en nod drar ned den, och nod B gör det.
\ans{b} Nod A och nod C. Båda sände recessivt och läser tillbaka dominant. Nod D sände ingenting
och har inget att jämföra med.
\ans{c} Att någon annan sänder samtidigt, och att den andra har företräde. Sker det i
arbitreringsfältet har A och C förlorat arbitreringen och blir mottagare.

\solution{ex:1:3}{Terminering}
\ans{a} Det sitter tre termineringsmotstånd i stället för två: tre parallella 120~\textohm{} ger
40~\textohm. Troligen har en nod mitt på bussen ett motstånd påslaget.
\ans{b} För låg terminering dämpar signalen men förstör den inte på en kort kabel vid måttlig
hastighet --- flankerna hinner ändå ställa sig. Felet visar sig när kabeln blir längre, farten
högre eller miljön störigare, alltså typiskt när konstruktionen flyttas från labbänken.

\solution{ex:1:4}{Asymmetrin}
\ans{a} Ja. Arbitreringen kräver bara att \emph{ett} av de två tillstånden vinner entydigt; vilket
spelar ingen roll för mekanismen.
\ans{b} Högst identifierare skulle vinna, eftersom en etta då skrev över en nolla.
\ans{c} Tänk på vad bussen gör när ingen sänder: den vilar i det tillstånd ingen driver, alltså
recessivt. Med den valda ordningen är viloläget det \emph{förlorande} tillståndet, vilket betyder
att en nod som slutar sända automatiskt drar sig ur. Vore det tvärtom skulle en tyst buss se ut
som en nod som vinner allt, och en nod som tappade strömmen skulle blockera bussen.

\section{Framen}

\solution{ex:2:1}{Läs framen}
\ans{a} $1 + 12 + 6 + 40 + 16 + 2 + 7 = 84$ bitar. SOF 1, arbitreringsfält 12 (11 identifierare +
RTR), kontrollfält 6, data $5 \cdot 8 = 40$, CRC-fält 16 (15 + avgränsare), ACK-fält 2, EOF 7.
\ans{b} 84~µs vid 1~Mbit/s, före stoppning. Med stoppbitar och de tre bitarna intermission blir
det i praktiken något mer.
\ans{c} \code{0x2A6}, eftersom lägre numeriskt värde vinner.

\solution{ex:2:2}{Ingen adress}
\ans{a} En. Sensorn sänder temperaturen en gång, med en identifierare som beskriver
\emph{innehållet}.
\ans{b} Två frames, en till vardera mottagaren, eller en broadcast-adress som protokollet måste
definiera särskilt.
\ans{c} Ingenting. Den tredje noden börjar bara lyssna efter samma identifierare. Det är CAN:s
starkaste argument: mottagare kan läggas till utan att sändaren vet om det.

\solution{ex:2:3}{Kvittensen}
\ans{a} Att minst en nod tog emot framen med korrekt CRC.
\ans{b} Hur många som gjorde det, vilka de var, om någon av dem brydde sig om innehållet, och om
någon faktiskt behandlade den. Kvittensen gäller överföringen, inte behandlingen.
\ans{c} Den läser recessivt, eftersom ingen annan finns som kan dra ned luckan. En standardenlig
kontroller tolkar det som ett kvittensfel, ökar sin sändarfelräknare och skickar om --- vilket är
varför en ensam nod på en CAN-buss sänder samma frame om och om igen tills någon svarar.

\solution{ex:2:4}{DLC}
\ans{a} Fältet ligger i kontrollfältet vars bredd var bestämd, och fyra bitar var vad som fanns.
Att bara 0--8 används är priset för ett fält som är jämnt i bitantal snarare än snålt.
\ans{b} Drivrutinens. Registerlagret vidarebefordrar värdet som det är, och kontrollern
sekvenserar efter det --- ingen av dem validerar. Att lägga kontrollen på båda ställena är precis
så de två kopiorna av regeln glider isär; den hör hemma där framen skapas.

\section{Bitstoppning och CRC-15}

\solution{ex:3:1}{Stoppa en ström}
\ans{a} \code{0000011111} blir \code{000001111101}: en stoppbit efter de fem nollorna och en efter
de fem ettorna. 10 bitar blir 12.
\ans{b} \code{1111111111} blir \code{111110111110}: en stoppbit efter de första fem ettorna, och
eftersom stoppbiten själv bryter serien räknas de följande fem ettorna om från början, vilket ger
ytterligare en. 10 bitar blir 12.
\ans{c} \code{0101010101} är oförändrad. Ingen serie når fem, så ingen stoppbit behövs. 10 bitar
förblir 10.

\solution{ex:3:2}{Avstoppa}
\ans{a} Biten på position 6 (nollan efter fem ettor) och biten på position 11 (ettan efter fem
nollor). Positionerna räknas från ett.
\ans{b} \code{1111100000}.
\ans{c} Ett stoppfel, efter den sjätte biten. Efter fem lika bitar \emph{måste} nästa vara den
motsatta; är den det inte är bitströmmen trasig, och mottagaren förkastar framen.

\solution{ex:3:3}{CRC-motorn}
Registret efter varje bit, med start från noll:

\begin{narrowtable}{@{}llll@{}}
\thead{Bit} & \thead{Värde} & \thead{Registret efter} & \thead{Hex} \\ \midrule
1 & 1 & \code{100010110011001} & \code{0x4599} \\
2 & 1 & \code{000101100110010} & \code{0x0B32} \\
3 & 0 & \code{001011001100100} & \code{0x1664} \\
4 & 1 & \code{110100101010001} & \code{0x6951} \\
\end{narrowtable}

Lägg märke till steg 2 och 3: när registrets översta bit och den inkommande biten är lika blir
återkopplingen noll, och steget är en ren skiftning utan XOR.

\solution{ex:3:4}{Varför noll?}
\ans{a} Checksumman är resten vid division med polynomet. Att lägga resten sist i dividenden gör
den jämnt delbar, och en jämn division lämnar resten noll. Mottagaren behöver därför ingen
jämförelse --- den tittar bara efter om registret är nollställt.
\ans{b} Samma modul kan användas på både sändar- och mottagarsidan, utan lägesomkoppling och utan
en separat jämförelsekrets. I hårdvara är det en påtaglig besparing, och en felkälla mindre.
\ans{c} Sändarens och mottagarens checksummor beräknas då över olika bitströmmar, eftersom
stoppbitarna beror på datan. Resultatet blir att varje frame som innehåller minst en stoppbit
förkastas som CRC-fel --- alltså nästan alla. Symptomet är förvirrande: korta frames med växlande
data fungerar, längre frames gör det inte.

\section{Arbitrering}

\solution{ex:4:1}{Två noder}
\ans{a} \code{0x155} är \code{00101010101} och \code{0x154} är \code{00101010100}.
\ans{b} Vid bit 11, den sista biten i identifieraren.
\ans{c} B vinner, eftersom den sänder dominant där A sänder recessivt. A upptäcker det omedelbart,
slutar sända, blir mottagare för B:s frame och försöker igen när bussen blivit ledig.

\solution{ex:4:2}{Tre noder}
Binärt: \code{0x400} är \code{10000000000}, \code{0x200} är \code{01000000000} och \code{0x201} är
\code{01000000001}.
\ans{a} Vid bit 1. A sänder recessivt medan B och C sänder dominant, så A faller bort direkt.
\ans{b} Vid bit 11. B och C är identiska i tio bitar och skiljer sig först i den sista.
\ans{c} B, med identifieraren \code{0x200}.

\solution{ex:4:3}{Samma identifierare}
\ans{a} Ingenting. Bitarna är identiska hela arbitreringsfältet igenom, så ingen av noderna ser
någon skillnad mellan sig själv och bussen. Båda tror att de vann.
\ans{b} Så fort datafälten skiljer sig sänder den ena noden recessivt och läser dominant. Utanför
arbitreringsfältet är det inte en förlorad arbitrering utan ett \emph{bitfel}: noden avbryter med
en felframe, båda frames förstörs, och båda noderna skickar om --- varpå samma sak händer igen.
\ans{c} Därför att det inte kan uppstå i drift av en korrekt konfigurerad buss. Identifierare
tilldelas per sändande nod vid systemdesignen; två noder med samma identifierare är ett fel i
tilldelningen, inte något som händer av sig självt.

\solution{ex:4:4}{Svält}
\ans{a} Den förlorar varje arbitrering mot varje annan nod, eftersom \code{0x7FF} är det högsta
möjliga värdet och därmed lägst prioritet. Är bussen alltid upptagen kommer den aldrig fram.
\ans{b} Nej. Protokollet garanterar att den \emph{högst} prioriterade framen kommer fram inom en
förutsägbar tid, och det är vad determinism i CAN betyder. Att lågprioriterad trafik kan svälta
är en konsekvens av prioritetsordningen, och den ordningen är poängen.
\ans{c} Till exempel: hålla bussbelastningen nere (30--50~\% är tumregeln i fordonssystem);
planera identifierarna efter vad som är tidskritiskt; periodisera sändningarna så att ingen nod
sänder så ofta den kan; eller dela upp trafiken på flera bussar.

\section{Bittajming}

\solution{ex:5:1}{Räkna tajmingen}
\ans{a} $50\,000\,000 / 1\,000\,000 = 50$ klockcykler per bit.
\ans{b} $50\,000\,000 / 500\,000 = 100$ klockcykler per bit.
\ans{c} $50 \cdot 0{,}7 = 35$, alltså vid tick 35 av 50.

\solution{ex:5:2}{Sampelpunkten}
\ans{a} För att signalen ska hinna ut till bussens bortersta nod och tillbaka innan biten läses.
Det är propagationssegmentets uppgift, och sampelpunkten måste ligga efter det.
\ans{b} Under arbitreringen skulle en nod läsa bussen innan en annan nods dominanta bit hunnit
fram, och alltså tro att den vann när den förlorade. Två noder skulle då sända samtidigt resten
av framen.
\ans{c} Propagationssegmentet. Det är det segment som representerar gångtiden, och det är därför
kabellängd och bithastighet hänger ihop.

\solution{ex:5:3}{Stoppning och synk}
\ans{a} Fem.
\ans{b} Bitstoppningen, femregeln. Efter fem lika bitar sätts alltid en bit med motsatt värde in,
vilket garanterar en flank.
\ans{c} Mottagarna skulle kunna gå långa sträckor utan flanker att resynkronisera mot, och driva
i väg tills de samplar fel bit. Det skulle i praktiken kräva en klockledning eller betydligt
noggrannare oscillatorer --- vilket är själva skälet till att stoppningen finns.

\solution{ex:5:4}{Två symptom}
\ans{a} Ingenting fungerar. Noderna är ur fas redan från SOF, så varje frame blir obegriplig i
andra änden. Felet syns omedelbart och är därför det snällare av de två.
\ans{b} De flesta bitar läses lika, men de som ligger nära en flank eller påverkas av
propagationsfördröjning läses olika. Symptomet blir sporadiska CRC- och stoppfel som blir värre
med längre kabel och högre belastning. Det är svårare att felsöka därför att det ser ut som en
dålig kabel eller en störig miljö, och därför att det fungerar tillräckligt bra för att man ska
tro att grundinställningen är rätt.

\section{Felhantering}

\solution{ex:6:1}{Vem ser felet?}
\ans{a} CRC-fel. Mottagarna upptäcker det; sändaren får veta det indirekt, genom att ingen
kvitterar.
\ans{b} Formatfel. Alla noder ser det, eftersom alla vet vilka bitar som måste vara recessiva.
\ans{c} Kvittensfel. Bara sändaren ser det: den släpper ACK-luckan och läser tillbaka recessivt.
\ans{d} Stoppfel. Alla noder ser det, eftersom alla räknar lika bitar i rad.

\solution{ex:6:2}{Räknarna}
\ans{a} $8 \cdot 8 = 64$.
\ans{b} $64 - 20 = 44$.
\ans{c} Error active i båda fallen, eftersom gränsen till error passive går vid 127.

Lägg märke till proportionerna: åtta fel kostar 64, och det krävs 64 lyckade sändningar för att
arbeta av dem. Räknaren är medvetet trögare på vägen ned än på vägen upp.

\solution{ex:6:3}{Asymmetrin}
\ans{a} Därför att räknarna inte mäter hur många fel som inträffat, utan hur sannolikt det är att
felet är nodens eget. En sändande nod är inblandad i varje fel den ser; en mottagande nod är en av
många som såg samma fel. Den som sänder är alltså en mer trolig felkälla, och straffas därefter.
\ans{b} Den trasiga noden är inblandad i \emph{varje} fel och klättrar åtta steg i taget, medan
den friska noden på en störig buss ökar ett steg i taget och hinner arbeta av sig mellan
störningarna. Efter en stund är den trasiga noden error passive eller bus off medan den friska
fortfarande är error active.

\solution{ex:6:4}{Bus off}
\ans{a} Därför att en nod med en trasig sändare rapporterar fel på varje frame --- också på de
frames som är helt korrekta. Dess felframes är dominanta och förstör andras trafik, så en enda
trasig nod kan lägga hela bussen.
\ans{b} Den hade blivit oanvändbar. Varje frame från varje nod hade avbrutits av den trasiga
nodens felflagga, och all trafik hade upphört.
\ans{c} Därför att felet kan vara övergående: en lös kontakt, en störning, en spänningsdipp. En
permanent avstängning hade krävt att någon åker ut och startar om noden för hand, också när felet
sedan länge är borta.

\section{Den förenklade kontrollern}

\solution{ex:7:1}{Vad saknas}
\ans{a} Fullständig: förloraren blir mottagare och tar emot den vinnande framen, och skickar om
sin egen automatiskt när bussen blivit ledig. Förenklad: förloraren sätter sin felflagga, lämnar
sändarrollen och går till vila --- den tar inte emot framen som vann, och skickar inte om.
\ans{b} Fullständig: mottagaren avbryter med en felframe, alla noder kastar framen, och sändaren
skickar om. Förenklad: mottagaren förkastar framen och sätter felflaggan. Ingen felframe sänds,
så sändaren vet ingenting.
\ans{c} Fullständig: sändaren upptäcker kvittensfelet och skickar om. Förenklad: sändaren läser
aldrig tillbaka ACK-luckan, så överföringen fullbordas precis som om allt gått bra.
\ans{d} Fullständig: båda hamnar i mottagningskön. Förenklad: den andra skriver över den första
om den första inte hunnit kvitteras, och den nyaste vinner.

\solution{ex:7:2}{Klibbiga bitar}
\ans{a} $40 \cdot 10^{-6} \cdot 50 \cdot 10^{6} = 2000$ klockcykler.
\ans{b} Så gott som ingenting. Pulsen är en cykel lång och sannolikheten att en pollning råkar
träffa just den cykeln är ungefär 1 på 2000. Mjukvaran skulle missa nästan varje händelse, och de
få den fångade skulle se slumpmässiga ut.
\ans{c} Därför att en automatisk nollställning vid läsning gör det omöjligt att läsa statusen två
gånger, och därmed omöjligt att skriva kod som först kontrollerar en bit och sedan agerar på den.
En uttrycklig nollställning skiljer dessutom "jag har sett händelsen" från "jag har hanterat
den" --- och det är mjukvaran, inte hårdvaran, som vet när det senare är sant.

\solution{ex:7:3}{Två steg}
\ans{a} Den förväxlar \emph{porten är fri igen} med \emph{framen kom fram}. Statusbiten kommer
tillbaka både efter en lyckad sändning och efter en förlorad arbitrering.
\ans{b} Vänta på att porten blir fri, och \emph{sedan} läsa felflaggan. Först den kombinationen
säger hur det gick.
\ans{c} Förlorad arbitrering, stoppfel, misslyckad CRC-kontroll, och en mottagen fjärrframe. Alla
fyra ger samma enda bit, så drivrutinen kan inte skilja dem åt.

\solution{ex:7:4}{Gränsen}
\ans{a} Till exempel genom att buffra mottagna frames i en kö, så att mjukvaran aldrig märkte att
kontrollern bara håller en. Eller genom att automatiskt nollställa statusbiten vid läsning, så
att mjukvaran slapp kvittera.
\ans{b} Ett dolt beteende är ett beteende man inte kan resonera om. En kö som ibland svämmar över
ger frames som försvinner utan spår, och felsökningen börjar på fel ställe: i mjukvaran, som
gjorde rätt. En begränsning man känner till går att planera för; en man inte känner till går
bara att stöta på.
\ans{c} Att ett gränssnitt ska vara \emph{ärligt} snarare än hjälpsamt. Ett lager som döljer
den underliggande hårdvarans egenskaper flyttar inte problemet, det gör det bara osynligt --- och
det som är osynligt går inte att testa, mäta eller resonera om. Samma princip gäller ett
drivrutins-API: det ska säga vad hårdvaran gör, inte vad man önskar att den gjorde.
